\documentclass[12pt,a4paper]{article}
\usepackage[utf8]{inputenc}
\usepackage[T1]{fontenc}
\usepackage[french]{babel}
\usepackage{geometry}
\usepackage{amsmath}
\usepackage{amssymb}
\usepackage{parskip}
\usepackage{titlesec}
\usepackage{xcolor}
\usepackage{booktabs}
\usepackage{enumitem}

\geometry{margin=2.5cm}

\definecolor{mainblue}{RGB}{30, 80, 160}
\definecolor{lightgray}{RGB}{245, 245, 245}

\titleformat{\section}{\large\bfseries\color{mainblue}}{}{0em}{}[\titlerule]
\titleformat{\subsection}{\normalsize\bfseries}{}{0em}{}

\title{
    \vspace{-1cm}
    {\large Marchés Électroniques — Projet 1}\\[0.4em]
    {\LARGE\bfseries\color{mainblue} Résumé du Projet}\\[0.3em]
    {\normalsize Modèle d'Almgren-Chriss : Stratégies de Liquidation Optimale}
}
\author{Pour Fernando — version de prise en main rapide}
\date{Avril 2026}

\begin{document}

\maketitle
\thispagestyle{empty}

% ─────────────────────────────────────────────
\section{C'est quoi le problème en deux phrases ?}

On a un gros paquet d'actions à vendre. On ne peut pas tout vendre d'un coup (ça ferait chuter le prix). On cherche \textbf{comment étaler les ventes dans le temps} pour minimiser les coûts tout en limitant le risque.

\medskip
\colorbox{lightgray}{\parbox{\linewidth}{\small
\textbf{Intuition clé :} Vendre vite $\Rightarrow$ impact de marché élevé (coût fort, risque faible).
Vendre lentement $\Rightarrow$ impact faible, mais le prix peut bouger pendant qu'on attend (risque fort).
Il faut trouver le bon équilibre.
}}

% ─────────────────────────────────────────────
\section{Notations de base}

\begin{itemize}[noitemsep]
    \item $Q$ : quantité initiale d'actions à vendre
    \item $T$ : horizon de temps (durée totale pour tout vendre)
    \item $q(t)$ : inventaire restant à l'instant $t$ (on part de $Q$, on finit à $0$)
    \item $v(t) = -\dot{q}(t)$ : vitesse de vente à l'instant $t$
    \item $\eta$ : coefficient d'impact de marché (plus $\eta$ est grand, plus vendre vite coûte cher)
    \item $\sigma$ : volatilité du prix
    \item $\lambda$ : aversion au risque (paramètre du trader)
    \item $\kappa = \sqrt{\lambda\sigma^2/\eta}$ : taux caractéristique qui règle la forme des stratégies
\end{itemize}

% ─────────────────────────────────────────────
\section{Section 1 — Trois types d'ordres classiques}

On compare trois stratégies optimales selon le \textbf{benchmark de prix} utilisé.

\subsection{Implementation Shortfall (IS) — Benchmark : prix initial $S_0$}

On mesure le coût par rapport au prix au moment où on \emph{décide} de vendre.

$\Rightarrow$ La stratégie optimale est \textbf{agressive au début} (on vend beaucoup au départ puis on ralentit).

\[
q^*(t) = Q \cdot \frac{\sinh(\kappa(T-t))}{\sinh(\kappa T)}
\]

Logique : on veut réduire l'inventaire rapidement pour ne pas être exposé trop longtemps à la volatilité.

\subsection{Target Close (TC) — Benchmark : prix final $S_T$}

On mesure le coût par rapport au prix de \emph{clôture}.

$\Rightarrow$ La stratégie optimale est \textbf{lente au début, puis accélère vers la fin}.

\[
q^*(t) = Q\left(1 - \frac{\sinh(\kappa t)}{\sinh(\kappa T)}\right)
\]

Logique : on attend le plus longtemps possible pour coller au prix final.

\subsection{TWAP — Benchmark : prix moyen sur $[0,T]$}

On mesure le coût par rapport au \emph{prix moyen} sur toute la période.

$\Rightarrow$ La stratégie optimale est toujours \textbf{linéaire}, indépendamment de $\lambda$ !

\[
q^*(t) = Q \cdot \frac{T-t}{T}
\]

Logique : vendre à rythme constant minimise simultanément coût et variance.

\medskip
\colorbox{lightgray}{\parbox{\linewidth}{\small
\textbf{Résultat remarquable :} IS et TC ont la même frontière efficace (symétrie temporelle). Le TWAP est un point unique sur cette frontière.
}}

% ─────────────────────────────────────────────
\section{Section 2 — Méthode numérique (Euler + Shooting)}

Le but ici : \textbf{retrouver numériquement} les solutions analytiques de la Section 1.

\subsection{Pourquoi faire ça ?}

Les formules de la Section 1 existent car le problème est simple. Mais dans des cas plus complexes (actifs multiples, contraintes...), on n'a pas de formule. Il faut des méthodes numériques.

\subsection{Ce qu'on fait concrètement}

\begin{enumerate}[noitemsep]
    \item L'optimisation donne une \textbf{équation différentielle} : $\ddot{q}(t) = \kappa^2 q(t)$
    \item On connaît $q(0) = Q$ mais pas $v(0) = \dot{q}(0)$ (la vitesse initiale)
    \item On \textbf{tire} (shooting) : on essaie plusieurs valeurs de $v(0)$ et on intègre numériquement avec le schéma d'Euler
    \item On cherche par bissection la valeur de $v(0)$ telle que $q(T) = 0$
\end{enumerate}

\medskip
\colorbox{lightgray}{\parbox{\linewidth}{\small
\textbf{Résultat :} La solution numérique coïncide avec la formule analytique — validation croisée réussie.
}}

% ─────────────────────────────────────────────
\section{Section 3 — Deux actifs avec couverture croisée}

On passe à un problème plus réaliste : \textbf{deux actifs corrélés}.

\subsection{Setup}

\begin{itemize}[noitemsep]
    \item $S^1$ : actif illiquide (fort impact $\eta_1 = 0.2$) — on \emph{doit} le vendre entièrement
    \item $S^2$ : actif liquide (faible impact $\eta_2 = 0.01$) — on peut aller short pour \emph{couvrir}
    \item Corrélation $\rho = 0.95$ : les deux actifs bougent presque ensemble
\end{itemize}

\subsection{Idée clé}

Quand on détient beaucoup de $S^1$, on est exposé à un risque de prix fort. En shortant $S^2$ (qui est très corrélé), on réduit cette exposition.

\[
\text{Coût total} = \underbrace{\eta_1 \int v_1^2\,dt + \eta_2 \int v_2^2\,dt}_{\text{coût d'impact}}
+ \lambda \underbrace{\int \left(\sigma_1^2 q_1^2 + 2\rho\sigma_1\sigma_2 q_1 q_2 + \sigma_2^2 q_2^2\right)dt}_{\text{variance du portefeuille}}
\]

Le terme croisé $2\rho\sigma_1\sigma_2 q_1 q_2$ devient \textbf{négatif} si $q_2 < 0$ (position short), ce qui réduit la variance.

\subsection{Méthode}

On utilise la \textbf{programmation dynamique de Bellman} sur une grille $(t, q_1, q_2)$ : on remonte en arrière dans le temps pour calculer la stratégie optimale.

\medskip
\colorbox{lightgray}{\parbox{\linewidth}{\small
\textbf{Résultat :} Shorter $S^2$ réduit la variance de \textasciitilde20-30\% pour un coût supplémentaire modeste — la couverture croisée est rentable.
}}

% ─────────────────────────────────────────────
\section{Section 4 (Bonus) — Réseaux de neurones}

\subsection{Objectif}

Apprendre la stratégie optimale \textbf{sans formule analytique}, juste en entraînant un réseau de neurones.

\subsection{Comment ça marche}

\begin{itemize}[noitemsep]
    \item On paramètre la politique : $v_t = \pi_\theta(t/T,\; q_t/Q)$ (un MLP avec poids $\theta$)
    \item La fonction de perte est le coût moyen-variance simulé par Monte Carlo
    \item On minimise avec Adam (gradient descent)
    \item Trois réseaux entraînés : IS, TC, et POV (participation au volume)
\end{itemize}

\medskip
\colorbox{lightgray}{\parbox{\linewidth}{\small
\textbf{Résultat :} Les réseaux IS et TC retrouvent les solutions analytiques. Le réseau POV apprend une stratégie sans formule connue — preuve que l'approche est générale.
}}

% ─────────────────────────────────────────────
\section{Vue d'ensemble — Ce que fait le code}

\begin{center}
\begin{tabular}{lll}
\toprule
\textbf{Fichier} & \textbf{Contenu} & \textbf{Rôle} \\
\midrule
\texttt{src/almgren\_chriss.py} & Formules analytiques & Stratégies IS, TC, TWAP \\
\texttt{src/bellman.py} & Prog. dynamique & Bellman 1 actif et 2 actifs \\
\texttt{main.ipynb} & Notebook principal & Calculs + graphiques \\
\texttt{Section 1 - Report.tex} & Rapport Section 1 & Dérivations mathématiques \\
\bottomrule
\end{tabular}
\end{center}

% ─────────────────────────────────────────────
\section{Schéma mental global}

\begin{center}
\begin{tabular}{ccc}
\textbf{Section 1} & \textbf{Section 2} & \textbf{Section 3--4} \\
Formules exactes & Méthodes numériques & Extensions \\
1 actif, 3 benchmarks & Euler + Shooting & 2 actifs / NN \\
$\downarrow$ & $\downarrow$ & $\downarrow$ \\
Solution analytique & Validation numérique & Généralisation \\
\end{tabular}
\end{center}

\vfill
\hrule
\smallskip
{\small\textit{Ce document est un résumé de prise en main. Pour les démonstrations complètes, voir \texttt{Section 1 - Report.tex} et \texttt{main.ipynb}.}}

\end{document}
